To separate computation security from physical assumptions, write $\mathsf{Acc}$ for gateway acceptance, $\mathsf{Cov}$ for the event that both the Byzantine and honest-miscoverage allowances hold, and $\mathsf{Comp}$ for the event that the report-bound true reference state lies in the authorized policy. For the chosen adversary resources and concrete parameters, let $\varepsilon_{\rm snark}$, $\varepsilon_{\rm sig}$, $\varepsilon_{\rm MiMC}$, $\varepsilon_{\rm ctx}$, and $\varepsilon_{\rm state}$ upper-bound failures, respectively, of the proof-system soundness guarantee, required signature unforgeability, signed-message prehash binding, deterministic context-map binding, and challenge-state enforcement.

\begin{proposition}[Composed conditional acceptance bound]\label{prop:composition}
Under trusted task configuration, roster, and gateway clock, and with the failure events defined above,
\begin{equation}\label{eq:composed}
\Pr[\mathsf{Acc}\land\mathsf{Cov}\land\neg\mathsf{Comp}]
\leq \varepsilon_{\rm snark}+\varepsilon_{\rm sig}
+\varepsilon_{\rm MiMC}+\varepsilon_{\rm ctx}+\varepsilon_{\rm state}.
\end{equation}
\end{proposition}
\begin{proof}
Suppose none of the stated cryptographic or state checks fails. An accepted proof then has a witness satisfying the configured task, report, roster, threshold, nonce, expiry, and signature constraints. Its Valid classification implies $\varnothing\ne\mathcal F_k\subseteq P$. Under $\mathsf{Cov}$, Proposition~\ref{prop:soundness} establishes $z^\star\in\mathcal F_k$, and hence $\mathsf{Comp}$. Thus occurrence of $\mathsf{Acc}\land\mathsf{Cov}\land\neg\mathsf{Comp}$ requires at least one of the enumerated failures. Applying the union bound yields~\eqref{eq:composed}.
\end{proof}

The bound concerns \emph{wrong acceptance while the physical coverage premise holds}; it is not the unconditional false-acceptance rate in deployment. An unconditional guarantee also requires estimating the probability and relevance of $\neg\mathsf{Cov}$. In the modeled single-process gateway without restart, atomic marking after successful verification prevents a second acceptance of the same challenge; the corresponding term for this specific replay behavior can be taken as $\varepsilon_{\rm state}=0$. No such zero-failure claim extends to process restarts, multiple gateways, or every possible state-management fault.

The prototype uses 256-bit challenge nonces, SHA-256 for report and context digests, Groth16 over BN254, and the BN254 twisted-Edwards EdDSA/MiMC implementation provided by gnark. This EdDSA instantiation is not Ed25519. The analysis composes stated assumptions but does not re-prove the cryptographic library or collapse its heterogeneous primitives into one overall security-bit estimate. Cryptography certifies authenticated evidence and its computation; physical coverage determines whether the result applies to the real event.

\subsection{Resistance to protocol attacks}
We first consider evidence tampering and context substitution. Changing the bounds of a signed region changes its message under~\eqref{eq:signedbox}, so, assuming signature security, it cannot pass in-circuit verification while retaining an honest issuer's original signature. Replacing the task, report, roster, threshold, challenge, or expiry changes the public statement reconstructed by the gateway. A proof made for a previous statement cannot simply be submitted for a different one without violating the relevant signature, digest-binding, or proof-system assumption. Missing evidence, malformed rosters, and incorrect decision codes lead to rejection or abort, not to a proved Inconsistent classification.

Challenge replay is checked by validating the nonce and expiry, verifying the proof, and atomically recording successful use. This prevents an already accepted nonce from authorizing another request within the same implemented gateway instance. The observed acceptance of one among eight concurrent replay copies is consistent with this guarantee. Failed verification does not consume an unexpired nonce so that legitimate retries remain possible. An attacker can therefore submit invalid proofs repeatedly to consume verification capacity; rate limits, resource quotas, persistent state, and distributed consistency must be supplied by the deployment. Nor does a fresh nonce automatically detect resubmission of the same report.

Finally, consider Byzantine sources that select their intervals strategically with knowledge of the public policy. Provided the Byzantine and honest-miscoverage budgets hold, Proposition~\ref{prop:soundness} gives $z^\star\in\mathcal F_k$ irrespective of malicious box placement. A true event outside $P$ therefore cannot yield a Valid feasible set wholly inside $P$, while a true event inside $P$ cannot yield an Invalid feasible set wholly outside it. Malicious sources may nevertheless reduce the frequency of decisive outcomes or withhold signatures and force abort. No corresponding guarantee follows if either budget or the honest coverage assumption fails. Our synthetic attacks use true-state and policy-boundary information to choose adversarial boxes, but they are not a substitute for adversarial evaluation in a deployed system.

The current protocol addresses evidence authentication, context integrity, and same-instance challenge replay. It does not ensure continuous issuer availability, resistance to denial of service, revocation, or dynamic operation with a partial roster; those protections require additional system mechanisms.

\subsection{Privacy boundary and information leakage}
Zero knowledge enables verification of the four-valued relation without exposing the issuer boxes and signatures that constitute the private witness. It does not hide all protocol information. In this implementation the gateway receives report bytes $R$ and combines them with the task and challenge fields to reconstruct public context $c$. Therefore $c$ is a deterministic binding value, not encryption of $R$ or a hiding commitment. If the report comes from a small enumerable space and the remaining context fields are available, an observer who sees $c$ may recompute candidate digests and perform an offline dictionary test. The gateway, which receives $R$ directly, cannot obtain report confidentiality from the proof's zero-knowledge property.

\begin{table}[htbp]\centering\footnotesize
\caption{Visibility and leakage boundaries for protocol participants.}\label{tab:privacy-boundary}
\begin{tabularx}{\linewidth}{@{}p{0.19\linewidth}p{0.37\linewidth}X@{}}\toprule
Role & Visible within the protocol model & Not concealed by the ZK guarantee\\\midrule
Evidence issuer & Its own observation and region, signing operation, and report-bound context & Its own evidence and any report data supplied by the host application\\
Vehicle/prover & Report, regions, and signatures required to build the witness & Secrecy of the witness from the prover itself is not an objective\\
Gateway/verifier & Report bytes, task/domain, policy, threshold, roster, nonce/expiry, public context, proof, classification & Report content, public policy, issuer identities, classification\\
Public transcript observer & Exposed statement, proof, classification, and transport metadata & Linkability of public fields and offline candidate tests for low-entropy reports\\\bottomrule
\end{tabularx}\end{table}

The four-valued classification itself reveals information. If only $q$ such outputs are considered, they admit at most $4^q$ different sequences and thus at most $2q$ bits of discrete output entropy. This counts the output alphabet only; it excludes information conveyed by other public fields and is not a location-privacy guarantee. If an attacker can adapt the spatial policy and observe subsequent classifications, the feedback may progressively reduce the candidate location region. The amount of leakage depends on policy geometry, query privileges, and prior information; no fixed localization accuracy or convergence rate follows merely from the $2q$ bound.

The prototype protects private witnesses only under an authorized, challenge-bound policy-query model. Unrestricted adaptive querying, network anonymity, hidden roster membership, and cross-business unlinkability are not provided. Applications requiring repeated location checks or stronger identity privacy must add external policy authorization and auditing, query-rate or budget controls, and suitable transport and identity protections.
